\section{Discussion}

\textbf{What the experiments collectively establish.} The internship began
by assuming input quality was the limit and found it was not (\F{F2}).
Surface proximity was then shown to be an unreliable proxy for anatomy
(\F{F1}); correspondence splits into identity and localisation (\F{F3});
initialisation matters for distal segments but not at the coxa (\F{F4});
and the objective has no explicit joint observation term (\F{F6}), so
anatomically different configurations can produce similar surfaces. Error is
structured by joint and specimen, not by seed (\F{F8}).

\textbf{Why this is a non-identifiability problem.} The scan supplies
surface geometry; the optimiser explicitly recovers shape, pose, scale and
deformation, and anatomical correspondence is an intended consequence of
those recovered parameters rather than an independently observed quantity.
Several latent configurations give geometrically similar surfaces, so the
surface observation does not uniquely determine the underlying anatomical
state: the inverse map is non-identifiable in the configurations tested here.
Each mechanism studied acts on that ambiguity: free deformation $\mathbf d$
can absorb a pose error (\S\ref{sec:deformation-freedom}); pointwise learned
votes can be locally right and globally inconsistent
(\S\ref{sec:learned-correspondence}), a mismatch between a pairwise
retrieval objective and a deployment problem that is really a globally
coupled assignment; and supervision at one place does not constrain another.

\textbf{Synthetic-to-real limits.} Learned components trained on synthetic
targets did not transfer at the trained perturbation regime
(\S\ref{sec:pointnet-init}), and the project's own synthetic reliability
estimate did not survive expert ground truth (\S\ref{sec:reliability}).

\textbf{Four distinctions.} Representability is not recoverability;
identity is not location; plausibility is not correspondence; fit quality is
not measurement validity. Each was the concrete reason a plausible
intervention failed. The central outcome is therefore not one configuration
but experimentally validated distinctions between geometric proximity,
anatomical correspondence, parameter recovery and measurement validity: a
diagnostic framework for SMILify and a design criterion for future
body-plan-agnostic models. A useful parametric representation is defined not
only by what it can represent, but by what anatomical meaning survives its
recovery from data. A reconstruction is not yet a phenotype.
